\subsubsection{Resultados esperados [Pendientes: 1. acordar umbrales de desempeño y contraste, 2. fijar cobertura mínima evaluable]}\label{sec:construccion-resultados}
El resultado de Construcción será una versión integrada E5 lista para evaluación con usuarios, acompañada de la biblioteca E4, el paquete reproducible E3 actualizado, la arquitectura E2 y los registros de E6 correspondientes a los incrementos. La condición de cierre será poder ejecutar el flujo completo bajo condiciones identificadas y comprobar, mediante pruebas vinculadas a requisitos, que cada componente intercambia e interpreta correctamente sus datos. Además, deberá existir un informe del piloto que permita reconocer las condiciones examinadas, las diferencias de aproximación y las restricciones restantes. Los criterios cuantitativos de cobertura, latencia y contraste se acordarán antes de evaluar su cumplimiento y no se reemplazarán por una apreciación visual de las animaciones.

OE1 quedará respaldado por consultas reproducibles y por la identificación de sus resultados, OE2 por los casos ejecutados y catalogados dentro de la cobertura acordada, y OE3 por configuraciones defendibles con evidencia de preparación y contraste. OE4 contará con la integración georreferenciada y las primeras comprobaciones de interacción. El cierre técnico deberá distinguir lo que funciona como prototipo exploratorio de lo que dispone de validación territorial, de manera que una demostración del software no se presente como habilitación para uso operativo. Si la agregación solamente resulta defendible bajo ciertas condiciones, esas condiciones deberán acompañar el producto evaluable.
